\section{Definition of Hodge ideals for \texorpdfstring{\mbox{\normalsize$\jepPubBold{Q}$}}{Q}-divisors}\label{definition}
In general, we obtain an upper bound for the terms in the filtration on $\jepPubScript{M}(h^{-\alpha})$
by restricting to the open subset where the support of ${\rm div}(h)$ is smooth, as follows.



\begin{proposition}\label{prop_for_def}
Given a nonzero $h\in\jepPubScript{O}_X(X)$ and a positive rational number $\alpha$, 
for every $k\geqslant 0$ we have 
$$F_k\jepPubScript{M}(h^{-\alpha})\subseteq \jepPubScript{O}_X\big((k+1)Z-\lceil D\rceil \big)h^{-\alpha},$$
where $D=\alpha\cdot {\rm div}(h)$ and $Z={\rm Supp}(D)$,
while $F_k\jepPubScript{M}(h^{-\alpha})=0$ for $k<0$.
\end{proposition}

\begin{proof}
Let $\iota\colon X_0\to X$ be an open immersion such that the codimension of its image in $X$ is $\geqslant 2$ and $Z\vert_{X_0}$ is smooth (though possibly disconnected).
Note that our constructions are compatible with restrictions to open subsets. Moreover, since $\jepPubScript{M}(h^{-\alpha})$ is clearly torsion-free, it follows that
$F_k:=F_k\jepPubScript{M}(h^{-\alpha})$ is torsion free, hence the canonical map $F_k\to \iota_*\big(F_k\vert_{X_0}\big)$ is injective. Therefore it is enough to prove the assertion on 
$X_0$, hence we may assume that $Z$ is smooth. However, in this case the assertion follows 
from Corollary~\ref{smooth_case}.
\end{proof}


We can now  define the Hodge ideals for $\jepPubBold{Q}$-divisors. Let $X$ be a smooth complex algebraic variety and $Z$ a reduced effective divisor on $X$.
Given an effective $\jepPubBold{Q}$-divisor $D$ with ${\rm Supp}(D)=Z$, we define coherent ideals sheaves $I_k(D)$  in $\jepPubScript{O}_X$ as follows.
Suppose first that there is a nonzero $h\in\jepPubScript{O}_X(X)$, with $H = {\rm div}(h)$, and a positive rational number $\alpha$ such that $D=\alpha H$. 
It turns out to be more convenient to work with the $\jepPubScript{D}_X$-module $\jepPubScript{M} (h^\beta)$, where $\beta = 1 - \alpha$. Recall that we have a 
filtered isomorphism 
$$\jepPubScript{M} (h^{-\alpha}) \longrightarrow \jepPubScript{M} (h^\beta), \,\,\,\,\,\,wh^{- \alpha} \longmapsto (wh^{-1})h^\beta,$$
and therefore, if $k\geqslant 0$,
it follows from Proposition~\ref{prop_for_def} that there is a unique coherent ideal $I_k(D)$ such that 
$$F_k\jepPubScript{M}(h^{\beta})=I_k(D)\otimes_{\jepPubScript{O}_X} \jepPubScript{O}_X\big(kZ + H \big)h^{\beta}$$
(note that we always have $\lceil D \rceil \geqslant Z$).
The definition is independent of the choice of $\alpha$ and $h$: indeed, using Remark~\ref{rem_behavior_etale}, it is enough to check this after the pullback by a suitable \'{e}tale surjective map, hence we deduce the independence assertion using Remark~\ref{rem_indep_equations1}. 
This implies that the general case of the definition follows by covering $X$ with suitable affine open subsets such that $D$ can be written as above in each of them. Note that when $D=Z$ we have $\beta = 0$, and so the ideals $I_k(D)$ are the Hodge ideals studied in \cite{MP1}.

\begin{remark}\label{rmk_inclusion_ideals}
From the definition and the filtration property,  it follows that we always have the inclusion 
$$\jepPubScript{O}_X (-Z) \cdot I_{k-1}(D) \subseteq I_k (D)\quad\text{for}\quad k\geqslant 1.$$ 
We note that for the reduced divisor $Z$,
we have the more subtle inclusions
$$I_k (Z) \subseteq I_{k-1} (Z)\quad \text{for}\quad k\geqslant 1$$
(see \cite[Prop.~13.1]{MP1}). We do not know however whether this holds for arbitrary $\jepPubBold{Q}$-divisors $D$, 
and in fact we suspect that this is not the case. (Note that it does hold when $D$ has simple normal crossings support
by Proposition \ref{Hodge_ideals_SNC}. It is also shown to hold when $D$ has an isolated weighted homogeneous singularity 
in the upcoming \cite{Zhang}.)
However, when $D = \alpha Z$ these inclusions do hold modulo the ideal $\jepPubScript{O}_X(-Z)$, see \cite[Cor.~B]{MP3}. 
More precisely, we have 
$$I_k(\alpha Z)+\jepPubScript{O}_X(-Z)\subseteq I_{k-1}(\alpha Z)+\jepPubScript{O}_X(-Z)\quad\text{for}\quad k\geqslant 1.$$
This implies in particular that if $I_k(\alpha Z)=\jepPubScript{O}_X$ for some $k\geqslant 1$, then $I_{k-1}(\alpha Z)=\jepPubScript{O}_X$.
\end{remark}

\begin{remark}\label{old_ideal}
According to Proposition \ref{prop_for_def}, we also have ideals $I_k^\prime (D)$ given by
$$F_k\jepPubScript{M}(h^{-\alpha})=I_k^\prime (D)\otimes_{\jepPubScript{O}_X} \jepPubScript{O}_X\big((k+1)Z -\lceil D\rceil \big)h^{-\alpha},$$
which are related to $I_k(D)$ by the formula
$$I_k (D) = I_k^\prime (D) \otimes_{\jepPubScript{O}_X} \jepPubScript{O}_X (Z -\lceil D\rceil).$$
\end{remark}

The following periodicity property often allows us to reduce our study to the case  $\lceil D\rceil = Z$.

\begin{lemma}\label{periodicity}
If $D'$ is an integral divisor with ${\rm Supp}(D')\subseteq {\rm Supp}(D)$, then
$$I_k(D+D')=I_k(D)\otimes_{\jepPubScript{O}_X} \jepPubScript{O}_X (- D').$$
In particular 
$$I_k (D) = I_k (B) \otimes \jepPubScript{O}_X (Z - \lceil D \rceil),$$
with $B = D + Z - \lceil D \rceil$ satisfying $\lceil B \rceil = Z$.
\end{lemma}
\begin{proof}
Using the notation in Remark \ref{old_ideal}, the equivalent statement
$$I_k^\prime (D+D')=I_k^\prime (D)$$
follows from the definition and Remark~\ref{rem_isom_comp_with_filtrations}.
\end{proof}

\begin{remark}\label{always_nontrivial}
Note that $I_k (D) \subseteq \jepPubScript{O}_X (Z - \lceil D \rceil)$ for all $k$, and so if $\lceil D \rceil \neq Z$, then one can never have 
$I_k (D) = \jepPubScript{O}_X$. It is however still interesting to ask whether $I_k (B) = \jepPubScript{O}_X$.
\end{remark}










